\section{The Bailout Protocol}

% ----------------------------------------------------------------
% 3.1  Overview
% ----------------------------------------------------------------

The Bailout Protocol is a deterministic, state-machine-driven procedure that
agents execute when their assigned task cannot be completed within a
predefined computation budget.  It exists to prevent indefinite resource
consumption, ensure that unresolved work is surfaced to a human supervisor,
and provide a structured escape path that never requires approval from any
subordinate machine actor.

The protocol is intentionally simple: an agent that encounters a bail-out
condition transitions through a fixed linear chain of states, each of which
represents a progressively higher level of human visibility, until a human
explicitly acknowledges the unresolved item.  No state may be skipped; no
transition may be reversed; no subordinate actor may invalidate or override
an active bail-out.

% ----------------------------------------------------------------
% 3.2  State Machine Definition
% ----------------------------------------------------------------

\subsection{State Space}

The Bailout State Machine $\mathcal{M}$ is a 6-tuple
\[
\mathcal{M} \;=\; \bigl\langle \Sigma,\;\Sigma_{\mathsf{initial}},\;
\Sigma_{\mathsf{terminal}},\;\rightarrow,\;L,\;T \bigr\rangle
\]
where:

\begin{itemize}
  \item $\Sigma = \{\,\text{IDLE},\; \text{BOUNDED},\; \text{BAIL\_PENDING},\;
         \text{ESCALATING},\; \text{HUMAN\_ACKNOWLEDGED},\; \text{RESOLVED}\,\}$
         is the finite set of states.
  \item $\Sigma_{\mathsf{initial}} = \{\,\text{IDLE}\,\}$ is the singleton
         set of initial states.
  \item $\Sigma_{\mathsf{terminal}} = \{\,\text{RESOLVED}\,\}$ is the
         singleton set of terminal states.
  \item $\rightarrow \;\subseteq\; \Sigma \times \Sigma$ is the deterministic
         transition relation defined in Table~\ref{tab:transitions}.
  \item $L : \Sigma \to \mathcal{P}(\text{Context})$ assigns to each state the
         set of context fields that must be preserved in the corresponding log
         entry.
  \item $T : \Sigma \to \mathbb{N}$ assigns to each state the maximum elapsed
         wall-clock seconds permitted before the protocol enforces the
         next transition automatically.
\end{itemize}

\bigskip
\noindent\textbf{State descriptions:}

\begin{enumerate}

  \item[\textbf{IDLE}] \label{state:IDLE}
        The agent is actively executing its primary task function.  No
        resource bounds have been exceeded and no bail-out procedure is
        active.  This is the \emph{normal operating state}.

  \item[\textbf{BOUNDED}] \label{state:BOUNDED}
        At least one resource bound (wall-clock time, token count, memory, or
        a custom metric) has been reached.  The agent is executing a
        structured mitigation routine --- typically a retry, a reduced-complexity
        fallback, or a partial-result compaction.  The agent still retains
        full autonomy; no human has been notified.

  \item[\textbf{BAIL\_PENDING}] \label{state:BAIL_PENDING}
        The BOUNDED mitigation routine has completed without resolving the
        task, or a bail-out trigger condition has fired asynchronously during
        BOUNDED.  The agent has entered the formal bail-out procedure.
        A human-in-the-loop notification is queued but not yet delivered.

  \item[\textbf{ESCALATING}] \label{state:ESCALATING}
        The human-in-the-loop notification has been dispatched and the
        protocol is waiting for an explicit acknowledgment signal.
        The agent may not continue task execution in this state.
        If no acknowledgment is received within the configured timeout
        $T(\text{ESCALATING})$, the machine performs an autonomous
        forced-transition to HUMAN\_ACKNOWLEDGED.

  \item[\textbf{HUMAN\_ACKNOWLEDGED}] \label{state:HUMAN_ACKNOWLEDGED}
        A human supervisor has confirmed receipt of the bail-out notification.
        The unresolved work item and all associated context are frozen.
        The agent is released to resume execution with modified parameters
        if the human has provided override instructions; otherwise it
        awaits further direction.

  \item[\textbf{RESOLVED}] \label{state:RESOLVED}
        The bail-out has been formally closed.  No further state transitions
        are possible.  The agent returns to IDLE.

\end{enumerate}

% ----------------------------------------------------------------
% 3.3  Transition Table
% ----------------------------------------------------------------

\subsection{Transition Relation}

Table~\ref{tab:transitions} enumerates every legal state transition,
its enabling condition, and the action performed atomically on entry to
the target state.

\begin{table}[htbp]
  \centering
  \caption{Bailout State Machine Transition Table}
  \label{tab:transitions}
  \begin{tabular}{@{}clll@{}}
    \toprule
    \textbf{\#} &
    \textbf{From} &
    \textbf{To} &
    \textbf{Enabling Condition} \\
    \midrule
    1 & IDLE       & BOUNDED       & Any resource bound $b \in B$ is reached                \\[4pt]
    2 & BOUNDED    & IDLE          & Mitigation routine returns \texttt{RESOLVED}           \\[4pt]
    3 & BOUNDED    & BAIL\_PENDING & Mitigation routine returns \texttt{UNRESOLVED} or
                                          timeout $T(\text{BOUNDED})$ expires \\[4pt]
    4 & BAIL\_PENDING & ESCALATING  & Entry action (notification dispatch) completes
                                          successfully                                               \\[4pt]
    5 & BAIL\_PENDING & IDLE        & Task becomes addressable via bypass mechanism
                                          (see \S\ref{sec:bypass})                                   \\[4pt]
    6 & ESCALATING & HUMAN\_ACKNOWLEDGED & Human acknowledgment signal received    \\[4pt]
    7 & ESCALATING & HUMAN\_ACKNOWLEDGED & Timeout $T(\text{ESCALATING})$ expires
                                          (autonomous forced transition)                             \\[4pt]
    8 & HUMAN\_ACKNOWLEDGED & RESOLVED & Human issues \texttt{CLOSE} directive         \\[4pt]
    9 & HUMAN\_ACKNOWLEDGED & IDLE    & Human provides override parameters; agent
                                          restarts task with modified budget                          \\[4pt]
    \bottomrule
  \end{tabular}
\end{table}

\bigskip
Transitions are deterministic in the sense that for any state
$\sigma \in \Sigma$, the set $\{\,\sigma' \mid (\sigma,\sigma') \in \rightarrow\,\}$
contains at most one element for any given triggering condition.
Multiple conditions may be evaluated simultaneously; the first satisfied
condition determines the transition.  When no condition is satisfied,
the machine remains in $\sigma$ (no implicit transition).

% ----------------------------------------------------------------
% 3.4  Bail-Out Trigger Condition
% ----------------------------------------------------------------

\subsection{Bail-Out Trigger Condition}

A bail-out is initiated whenever the following predicate evaluates to
\texttt{TRUE}:

\begin{definition}[Bail-Out Trigger]\label{def:bailout-trigger}
Let $B$ be the set of resource bounds configured for agent $a$, and let
$e_b(t)$ denote the current consumption of resource $b$ at wall-clock time $t$.
Define
\[
\phi_{\text{BOUNDED}}(a,t) \;=\; \bigvee_{b \in B}\; \text{bound\_exceeded}(b, e_b(t))
\]
to mean that \emph{at least one} configured bound has been breached.

The agent enters BOUNDED when $\phi_{\text{BOUNDED}}(a,t)$ becomes true.

The agent enters BAIL\_PENDING when it was in BOUNDED and either:

\begin{enumerate}
  \item $\phi_{\text{BOUNDED}}(a,t)$ remains true and the BOUNDED mitigation
        routine returns \texttt{UNRESOLVED}; \emph{or}
  \item the elapsed time in BOUNDED exceeds $T(\text{BOUNDED})$.
\end{enumerate}
\end{definition}

In practical terms, the trigger fires when the agent has exercised its
first-line mitigation (BOUNDED) and that mitigation has either failed
or timed out.

\bigskip
\noindent\textbf{Formal state invariant.}

At all times the machine satisfies:
\[
\sigma \in \Sigma \quad\Longleftrightarrow\quad
\begin{cases}
\text{active\_task}(a) \neq \varnothing          & \text{if } \sigma = \text{IDLE}   \\
\text{active\_task}(a) \neq \varnothing          & \text{if } \sigma = \text{BOUNDED} \\
\text{bailout\_context}(a) \neq \varnothing      & \text{if } \sigma \in
  \{\,\text{BAIL\_PENDING},\; \text{ESCALATING},\;
   \text{HUMAN\_ACKNOWLEDGED}\,\}                \\
\text{active\_task}(a) = \varnothing              & \text{if } \sigma = \text{RESOLVED}
\end{cases}
\]

% ----------------------------------------------------------------
% 3.5  Escalation Algorithm
% ----------------------------------------------------------------

\subsection{Escalation Algorithm}

When the machine enters BAIL\_PENDING, the following escalation procedure
is executed atomically before the transition to ESCALATING completes.

\begin{algorithm}[H]
\caption{Escalation Procedure}
\label{alg:escalation}
\begin{algorithmic}[1]
\Require $\text{agent } a$, $\text{task } \tau$, $\text{context } \kappa$,
         $\text{notification\_channel } \iota$
\Ensure $\text{notification sent}$, $\sigma \gets \text{ESCALATING}$

\State $\text{payload} \gets \text{serialize\_bailout\_report}(a, \tau, \kappa)$
\Comment{Freeze all relevant context into a structured log entry}
\State $\text{payload.bound\_history} \gets \text{read\_bound\_events}(a)$
\State $\text{payload.mitigation\_attempts} \gets
        \text{count\_mitigations}(a)$
\State $\text{payload.node\_id} \gets \text{agent.metadata.node\_id}$
\State $\text{payload.timestamp} \gets \text{utc\_now}()$

\State $\text{queue\_depth} \gets \text{query\_escalation\_queue}(\iota)$
\If{$\text{queue\_depth} > \text{MAX\_QUEUE\_DEPTH}$}
    \State $\text{fallback\_channel} \gets \text{get\_fallback\_channel}(\iota)$
    \State $\iota \gets \text{fallback\_channel}$
    \Comment{PagerDuty fallback, secondary email, SMS, etc.}
\EndIf

\State $\text{dispatch\_notification}(\iota,\; \text{payload})$
\State $\text{log\_escalation}(\text{payload})$
\Comment{Persist to audit trail before state transition commits}

\State $\sigma \gets \text{ESCALATING}$
\State $\text{start\_timer}(\text{ESCALATING},\; T(\text{ESCALATING}))$
\State \Return $(\sigma, \text{payload.notification\_id})$
\end{algorithmic}
\end{algorithm}

\bigskip
The escalation report payload includes the bound-exceeded history, the
mitigation attempts performed in BOUNDED, the originating agent's node
identifier, and a UTC timestamp.  This ensures a human reviewer has a
complete picture of the failure mode without needing to query external
systems.

\bigskip
\noindent\textbf{Escalation levels.}

If the acknowledgment is not received within $T(\text{ESCALATING})$,
the algorithm performs an autonomous forced transition to
HUMAN\_ACKNOWLEDGED (transition 7, Table~\ref{tab:transitions}).
This is a deliberate design choice: silence after a configured window
is treated as human acknowledgment by necessity, because any alternative
--- looping indefinitely or dropping the item --- is strictly worse.

% ----------------------------------------------------------------
% 3.6  Bypass Mechanism
% ----------------------------------------------------------------

\subsection{Bypass Mechanism}
\label{sec:bypass}

The bypass mechanism is a deterministic escape hatch that allows an agent
to exit the Bailout Protocol and return to IDLE \emph{without} requiring a
human acknowledgment signal.  Its existence is justified by a single
operational principle:

\begin{principle}[Bypass Justification]
A subordinate machine actor that can prove to the protocol that it has
identified a provably correct resolution to the original task constraint
may self-release from the bail-out procedure, because no human judgment
is required to verify the correction.
\end{principle}

\bigskip
The bypass is \emph{not} a shortcut.  It may only be invoked when the
agent can produce a verifiable proof object $P$ satisfying the following
formal predicate:

\begin{definition}[Bypass Predicate]\label{def:bypass-predicate}
An agent $a$ holding task $\tau$ in state $\sigma = \text{BAIL\_PENDING}$
may transition to IDLE (transition 5) if and only if there exists a
proof object $P$ such that:
\begin{align}
  \text{provably\_correct}(P, \tau) &= \texttt{TRUE} \label{eq:bypass-correct}\\
  \text{no\_human\_dependency}(P)   &= \texttt{TRUE} \label{eq:bypass-no-human}\\
  \text{bound\_reset}(P)            &= \texttt{TRUE} \label{eq:bypass-reset}
\end{align}
where:
\begin{itemize}
  \item \eqref{eq:bypass-correct} means that $P$ demonstrates $\tau$ satisfies
        all acceptance criteria with formal verification or
        deterministic re-execution;
  \item \eqref{eq:bypass-no-human} means that $P$ does not require any
        human-provided parameter, label, or judgment to be validated;
  \item \eqref{eq:bypass-reset} means that all resource bounds $b \in B$
        have been reset to their initial values so that the agent can
        re-enter IDLE with a clean budget.
\end{itemize}
\end{definition}

\bigskip
\noindent\textbf{Concrete examples of valid bypass scenarios:}

\begin{itemize}
  \item \emph{Cache-miss failure}: The agent discovers that the original
        failure was caused by a stale cached result; re-fetching the data
        from the canonical source resolves the task deterministically.
        The proof object $P$ is the fresh response from the canonical source.

  \item \emph{Token-bound retry}: The agent exceeded the token bound on a
        first attempt but has now loaded a smaller model variant that
        satisfies the same functional requirements within the budget.

  \item \emph{Pregame validation}: The agent identifies that the task
        request itself violates a precondition (e.g., requesting a file
        write to a path that does not exist).  Returning
        \texttt{INVALID\_REQUEST} to the caller satisfies the acceptance
        criteria without human intervention.
\end{itemize}

\bigskip
\noindent\textbf{Invalid bypass scenarios (bypass must not fire):}

\begin{itemize}
  \item Agent wants to extend the token bound and continue --- this is
        a human decision (requires HUMAN\_ACKNOWLEDGED);
  \item Agent encountered a transient network error and requests a retry
        with no change to the original parameters --- must go through
        the full escalation path;
  \item Agent cannot prove correctness but believes it is likely correct ---
        belief is insufficient; $P$ must be formal.
\end{itemize}

\bigskip
The bypass mechanism bypasses all machine actors by construction:
only the initiating agent's own proof object $P$ is evaluated against
Definition~\ref{def:bypass-predicate}.  No peer agent, no orchestration
layer, and no policy engine participates in the decision.  This ensures
that the bypass is an \emph{intra-agent} operation that cannot be blocked
by any external machine actor.

% ----------------------------------------------------------------
% 3.7  Timeout Handling
% ----------------------------------------------------------------

\subsection{Timeout Handling}

Each non-terminal state $\sigma \in \Sigma \setminus \{\text{RESOLVED}\}$
is assigned a timeout value $T(\sigma) \in \mathbb{N}$ (wall-clock seconds).
Timeouts are configured at the agent level and inherit from the
agent's policy profile unless explicitly overridden.

\bigskip
\noindent\textbf{Timeout semantics.}

When a timeout $T(\sigma)$ expires while the machine is in state $\sigma$,
the following handler is invoked:

\begin{algorithm}[H]
\caption{Timeout Handler}
\label{alg:timeout-handler}
\begin{algorithmic}[1]
\Require $\text{agent } a$, $\text{current state } \sigma$,
         $\text{elapsed } t_{\text{elapsed}}$, $\text{timeout } T(\sigma)$
\Ensure $\sigma'$ is the next state per Table~\ref{tab:transitions}

\State $\text{log\_timeout}(a,\; \sigma,\; t_{\text{elapsed}},\; T(\sigma))$
\Comment{Append to audit trail with full context}

\Switch{$\sigma$}
  \Case{$\text{BOUNDED}$}
    \State $\text{transition\_to}(\text{BAIL\_PENDING})$
    \Comment{Transition 3, second limb: mitigation timeout}
  \EndCase
  \Case{$\text{ESCALATING}$}
    \State $\text{transition\_to}(\text{HUMAN\_ACKNOWLEDGED})$
    \Comment{Transition 7: autonomous forced acknowledgment}
    \State $\text{notify\_audit\_log}(
           \text{"AUTONOMOUS\_ACK\_TIMEOUT"},\; a)$
  \EndCase
  \Case{$\text{BAIL\_PENDING}$}
    \State $\text{dispatch\_notification}(\iota,\; \text{payload})$
    \State $\text{transition\_to}(\text{ESCALATING})$
    \Comment{BOUNDED was never entered; direct escalation from IDLE timeout
             is not possible by design}
  \EndCase
  \Case{$\text{HUMAN\_ACKNOWLEDGED}$}
    \State $\text{log\_warning}(
           \text{"SUPERVISOR\_TIMEOUT"},\; a,\; \sigma)$
    \State $\text{send\_escalation\_pulse}(\text{fallback\_channel})$
    \Comment{Notify secondary supervisors; do not change $\sigma$}
  \EndCase
  \Default
    \State $\text{log\_error}(
           \text{"UNHANDLED\_TIMEOUT"},\; a,\; \sigma)$
    \State $\text{transition\_to}(\sigma)$
    \Comment{No-op; remain in current state and alert operator}
  \EndSwitch
\end{algorithmic}
\end{algorithm}

\bigskip
Timeouts are \emph{wall-clock based} and are evaluated by the agent's
local scheduler, not by an external orchestrator.  This ensures that
timeout handling is autonomous and does not introduce a distributed-systems
clock-synchronization dependency.

\bigskip
\noindent\textbf{Timeout configuration defaults:}

\begin{itemize}
  \item $T(\text{BOUNDED}) \;\;$    = 120 seconds (configurable per agent profile)
  \item $T(\text{ESCALATING}) \;$   = 300 seconds (5 minutes; human-facing channels)
  \item $T(\text{BAIL\_PENDING}) $  = 0 seconds (transition is instantaneous;
        no dwell time in BAIL\_PENDING by design)
  \item $T(\text{HUMAN\_ACKNOWLEDGED})$ is unbounded (a human may take
        as long as required to provide direction)
\end{itemize}

\bigskip
The BOUNDED timeout value (120 s) accommodates retry-and-fallback
mitigation routines that require I/O (file reads, API calls) while still
ensuring that a truly stalled agent does not remain in BOUNDED for an
excessive period.  The ESCALATING timeout (300 s) reflects the expected
human response latency for an async notification delivered to a
non-critical channel (e.g., email or Slack).  For critical-path integrations
(PagerDuty, SMS), a shorter $T(\text{ESCALATING})$ should be configured
to match the SLA expectation of the receiving human.

% ----------------------------------------------------------------
% 3.8  Context Preservation and Audit Log
% ----------------------------------------------------------------

\subsection{Context Preservation}

Each state transition emits a structured log entry to the agent's
audit trail.  The log entry schema is:

\begin{verbatim}
LogEntry {
  event_id:      UUID v4
  agent_id:      string
  task_id:       string
  from_state:    State
  to_state:      State
  trigger:       TransitionTrigger
  timestamp:     ISO 8601 UTC
  bound_state:   map[string, ResourceSnapshot]
  mitigation:    MitigationSummary | null
  bypass_proof:  ProofObject | null
  human_input:   HumanDirective | null
}
\end{verbatim}

\bigskip
The presence of a non-null \texttt{bypass\_proof} field on a transition
from BAIL\_PENDING to IDLE (transition 5) provides independent
verifiability that the bypass was exercised in accordance with
Definition~\ref{def:bypass-predicate}.  This allows post-hoc auditing
without requiring the human supervisor to be online at the time of the
bypass.

% ----------------------------------------------------------------
% 3.9  State Machine Summary
% ----------------------------------------------------------------

\subsection{Formal Summary}

Integrating the definitions above, the complete Bailout Protocol satisfies
the following properties:

\begin{enumerate}
  \item[\textbf{Determinism}] For any $(\sigma, c)$ pair, the transition
        relation $\rightarrow$ yields at most one target state.

  \item[\textbf{Progress}] Every non-terminal state has at least one
        outgoing transition; the machine cannot dead-lock in any state
        $\sigma \in \Sigma \setminus \{\text{RESOLVED}\}$.

  \item[\textbf{HumanOverride}] No agent may transition into RESOLVED
        without either a human \texttt{CLOSE} directive or an autonomous
        timeout-forced acknowledgment.

  \item[\textbf{BypassIsolation}] The bypass predicate (Definition 2)
        is evaluated entirely within the agent and is not subject to
        external veto by any machine actor.

  \item[\textbf{AuditCompleteness}] Every transition emits a
        \texttt{LogEntry} to the audit trail before the transition commits.
\end{enumerate}

\bigskip
The protocol balances autonomy (agents resolve bounded failures without
human intervention) with safety (unresolved failures always surface to a
human and cannot be suppressed by any subordinate actor).